\subsubsubsection{子集遍历}

\begin{lstlisting}[language=C++, basicstyle=\ttfamily\small]
for (int i = 0; i < 1 << n; i ++ )
\end{lstlisting}

这个代码可以从 $0$ 开始不重复的遍历每一个二进制数，因此时间复杂度 $O(2^n)$。

\subsubsubsection{子集枚举}

有的时候，我们希望枚举集合 $i$ 所对应的所有子集，显然我们不可能循环套循环求与的方法做，时间复杂度是 $O(2^{2n})$ 的，不优。

我们下面介绍子集枚举，他可以按照降序排序枚举集合。

\begin{lstlisting}[language=C++, basicstyle=\ttfamily\small]
for (int S = 0; S < 1 << n; S ++ )
    for (int T = S; T; T = (T - 1) & S)
\end{lstlisting}

$T$ 最初恰好是集合 $S$，在下一次循环时，我们让 $T = (T - 1) \& S$，容易发现，把 $T$ 的最后一个位掩码去掉了。

\begin{verbatim}
S = 1010100
T = 1010100 -> T = 1010011 & 1010100 = 1010000
\end{verbatim}

再进行下一次循环，我们把刚才的比特恢复，消去下一个比特。

\begin{verbatim}
S = 1010100
T = 1010000 -> T = 1001111 & 1010100 = 1000100
\end{verbatim}

接着再移去最后一个比特。

\begin{verbatim}
S = 1010100
T = 1000100 -> T = 1000011 & 1010100 = 1000000
\end{verbatim}

以此类推，不难发现，`T = (T - 1) \& S` 等价于我们倒序枚举连续集合的 `T -- `，相当于无视了中间的 $0$ 进行转移，这是优秀的，保证了每个子集唯一出现一次。

考虑每个枚举集合 $S$ 的子集 $T$ 的次数：

\[
\begin{aligned}
\sum_{S \subseteq \{1, 2, \cdots, n\}}\sum_{T \subseteq S}1 &= \sum_{i = 0}^n\binom{n}{i}2^i \\
&= (2 + 1)^n \\
&= 3^n
\end{aligned}
\]

因此子集枚举的复杂度是 $O(3^n)$。